
The lemma follows from \SGARefLink{sga3:xvii:lemma:3:9:ter}{3.9 ter} when the characteristic \(p\) of
\(k\) is zero, and from \SGARefLink{sga3:xvii:lemma:2:3:bis}{2.3 bis} when \(r=1\).  In the general case
\((p>0)\), realize \(G\) as an algebraic subgroup of
\((\mathbb G_a)^n\), for a suitable integer \(n\) (\SGARefLink{sga3:xvii:lemma:4:1:2}{4.1.2}), and
argue by induction on \(n-r\).  If \(r=n\), then indeed
\(G=(\mathbb G_a)^r\).  Otherwise, the quotient
\((\mathbb G_a)^n/G\) is a nonzero smooth, connected, unipotent
group.  By \SGARefLink{sga3:xvii:proposition:4-1-1}{4.1.1}, \SGARefLink{sga3:xvii:proposition:4:1:1:item:i}{(i)}\(\Rightarrow\)\SGARefLink{sga3:xvii:proposition:4:1:1:item:ii}{(ii)}, and \SGARefLink{sga3:xvii:lemma:2:3:bis}{2.3 bis}, it has a
composition series whose quotients are isomorphic to
\(\mathbb G_a\).  It follows that there is a smooth, connected
algebraic subgroup \(G_1\) of \((\mathbb G_a)^n\), containing \(G\),
such that
\[
  (\mathbb G_a)^n/G_1=\mathbb G_a.
\]
By induction, it suffices to show that \(G_1\) is isomorphic to
\((\mathbb G_a)^{n-1}\).

It is immediate that a homomorphism
\[
  \operatorname{Spec}k[X_1,\ldots,X_n]=(\mathbb G_a)^n
  \longrightarrow
  \mathbb G_a=\operatorname{Spec}k[T]
\]
is defined by an additive polynomial of the form
\[
  \sum_{i,j}a_{i,j}X_i^{p^j}.
\]
Since \(G_1\) is smooth, the linear part of this polynomial is
nonzero.  After a linear change of the coordinates \(X_i\), we may
assume that \(G_1\) is the algebraic subgroup of
\((\mathbb G_a)^n\) defined by
\[
  \tag{\(*\)}
  P(X)=-X_1+\sum_{j=1}^{q}a_jX_1^{p^j}
    +Q(X_2,\ldots,X_n)=0,
\]
\SGARefTarget{sga3:xvii:lemma:4:1:5:eq:star}
where
\[
  Q(X_2,\ldots,X_n)
    =\sum_{\substack{i>1\\j>0}}b_{i,j}X_i^{p^j}.
\]

We now argue by induction on the degree of \(P\).  If
\(\deg P=1\), it is clear that
\[
  G_1\xrightarrow{\ \sim\ }(\mathbb G_a)^{n-1}.
\]
Otherwise, since \(k\) is perfect, one has
\(Q(X)=Q_1(X)^p\), and we may define an endomorphism \(v\) of
\((\mathbb G_a)^n\) by
\[
\begin{aligned}
  X_i&\longmapsto X_i &&\text{for \(i>1\)},\\
  X_1&\longmapsto
    \sum_{j=1}^{q}a_j^{1/p}X_1^{p^{j-1}}+Q_1(X).
\end{aligned}
\]
\SGARefTarget{sga3:xvii:lemma:4:1:5:map:v}
It is clear that \(v\) induces an isomorphism on \(G_1\), and that
\(v(G_1)\) is defined in \((\mathbb G_a)^n\) by the equation
\[
  \tag{\(*\)\(_1\)}
  P_1(X)=-X_1+\sum_{j=1}^{q}a_j^{1/p}X_1^{p^j}
    +Q_1(X_2,\ldots,X_n)=0.
\]
\SGARefTarget{sga3:xvii:lemma:4:1:5:eq:star:1}
